\chapter{Conclusions}

This project successfully developed, modelled, and evaluated a position-holding joint mechanism for a to mitigate operational drift and enhance kinematic stability.

The Denavit-Hartenberg (D-H) kinematic model derived from the FaroArm QuantumX architecture was successfully integrated into a MATLAB-based error-propagation framework. Spatial uncertainty mapping confirmed that position uncertainty is heavily posture-dependent, ranging from a minimum $U_{95}$ spatial uncertainty of $56~\mu\text{m}$ near the base to a maximum of $74~\mu\text{m}$ at full extension. This mathematical model proved capable of evaluating single-joint locking scenarios to inform future multi-DOF arm designs.
	
Based on the error-propagation analysis, TPMs were established, including a static holding torque target of $7.725~\text{N}\cdot\text{m}$ (TPM\_03) and an operational articulation sweep of $90^\circ$ (TPM\_01). A qualitative evaluation chose an internal electromagnetic brake configuration as the optimal concept, as it provides a uniform clamping force that minimise parasitic deflections while maintaining full articulation range.

The full-scale mechatronic joint design successfully integrated mechanical, electronic, and software domains, specifying an IRF640N MOSFET driven by a 2N3904 NPN pre-amplifier stage to safely switch the 24~V, 8~A parallel coil load from an STM32 Nucleo microcontroller. Real-time angular position tracking would have been achieved via hardware timer encoder interfaces operating at 4000 counts per revolution. For the as-built electromagnet subassembly implementation, a low-side Darlington driver stage (pairing 2N3904 and TIP41C power transistors) was designed, simulated in LTspice, and replicated in parallel to safely validate high-current coil actuation during testing.
	
Physical testing of the built electromagnet verified the switching logic, current-sharing among parallel Darlington driver stages, and software safety protocols under operational currents up to $4.05~\text{A}$. Initial mechanical clamping was demonstrated. Although full $3~\text{mm}$ air-gap closure and strain-gauge force readings were constrained by the high bending stiffness of the mild steel pressure tabs, tight solenoid winding tolerances, and thermal limits of the fine wire gauge ($0.28~\text{mm}$ diameter), the experimental apparatus confirmed the overall feasibility of the electromagnetic locking concept.

Overall, the project established a theoretical foundation and partially functional subsystem prototype for position-holding PCMA joints. The generalized D-H error model and standardized design documentation provide a clear framework for future researchers to optimize actuator dimensions and material selection for full-scale arm integration.